\chapter{Writing and Maintaining Iceberg}
\label{ch:iceberg-write}

\begin{chapterintro}
Copy-on-write rewrites whole files on every update: slow writes, fast reads.
Merge-on-read writes small delete files instead: fast writes, slower reads until
compaction catches up. The choice is a function of your update rate and your read
latency budget, and it commits you to a maintenance schedule --- compaction,
snapshot expiration, orphan-file cleanup --- in a specific order. This chapter
covers the write modes, optimizing \texttt{MERGE}, and running maintenance
without deleting data a reader still needs.
\end{chapterintro}

\section{Copy-on-write versus merge-on-read}
\tbd

\section{\texttt{MERGE}, \texttt{UPDATE}, and \texttt{DELETE}}
\tbd

\section{Write distribution and skew on write}
\tbd

\section{Compacting data files}
\tbd

\section{Rewriting manifests}
\tbd

\section{Snapshot expiration and orphan files}
\tbd

\section{Metadata management for streaming tables}
\tbd

\section{Summary}
\tbd

\exercisehead
\begin{exercises}
\item Run a \texttt{MERGE} with and without the partition column in the
      \texttt{ON} clause. Explain the difference in files scanned from the plan.
\end{exercises}
